\subsection{\bt{Quality and valid solutions}{质量与有效方案}}
\label{sec:quality-eval}
\begin{figure*}[t]
\centering
\ifbilingual
\includegraphics[width=\textwidth]{figures/official-quality-zh}
\else
\includegraphics[width=\textwidth]{figures/official-quality}
\fi
\caption{\bt{Test quality and valid solutions in the official-API comparison. Each point uses all three cutoffs; unsuccessful sessions retain the entity-reference score. Positive values favor the method over no memory: $100(\mathrm{AUC}-\mathrm{AUC}_0)$ in percentage points, or $100(\mathrm{MAE}_0-\mathrm{MAE})/\mathrm{MAE}_0$ in percent. Shading marks the predeclared quality margin; panels have separate scales. Valid counts solutions passing the temporal audit and final evaluation, out of three attempts per task. The 90 sessions are one sequence per method and database, not independent repetitions.}{官方API对照的测试质量与有效方案。每点使用全部三个截点；失败会话保留实体表参照分数。正值表示优于无记忆：AUC为$100(\mathrm{AUC}-\mathrm{AUC}_0)$个百分点，MAE为$100(\mathrm{MAE}_0-\mathrm{MAE})/\mathrm{MAE}_0$百分比。阴影标出预设质量门槛；各子图使用独立刻度。有效方案数指每任务三次尝试中通过时间审计及最终评估的数量。90个会话对应每种方法、每个数据库一条序列，并非独立重复。}}
\label{fig:official-quality}
\Description{Six dot plots compare all five methods on the six RelBench tasks. Each panel shows quality relative to no memory, a zero line, the quality margin, and the valid-solution count. Unfavorable outcomes and fallback scores remain visible.}
\end{figure*}

\bi{Figure~\ref{fig:official-quality} pairs test quality with the number of usable solutions, using all 90 official-API sessions. Across all cutoffs, \system\ meets the predeclared margin on \OfFullQualityCount/6 tasks: $\mathrm{AUC}\geq\mathrm{AUC}_0-0.01$ or $\mathrm{MAE}\leq1.02\,\mathrm{MAE}_0$. Its post-votes MAE is 2.2\% lower than no memory, while user-badge and user-engagement AUC differ by +0.05 and -0.15 percentage points. All nine rel-stack attempts yield valid solutions, compared with eight for no memory, AWM, and G-Memory. These outcomes accompany the lower cost in Figure~\ref{fig:official-cumulative}; they do not imply that \system\ maximizes predictive quality on every task.}
{图~\ref{fig:official-quality}将测试质量与可用方案数量并列，覆盖全部90个官方API会话。全部截点上，\system\ 在\OfFullQualityCount/6项任务满足预设门槛：$\mathrm{AUC}\geq\mathrm{AUC}_0-0.01$或$\mathrm{MAE}\leq1.02\,\mathrm{MAE}_0$。post-votes MAE比无记忆低2.2\%；user-badge和user-engagement AUC分别相差+0.05与-0.15个百分点。rel-stack的九次尝试均得到有效方案，无记忆、AWM和G-Memory各为八次。这些结果与图~\ref{fig:official-cumulative}的较低代价相对应，但不意味着\system\ 在每项任务都取得最高预测质量。}

\bi{The rel-trial results expose the remaining failures. \system\ produces six valid solutions out of nine attempts, versus seven for ACE, AWM, and G-Memory. Its study-adverse MAE is \OfAdverseRegression\% higher than no memory, outside the margin; its first-cutoff attempt has no valid solution. On site-success, its mean improves but only one of three attempts is valid, so the mean alone would overstate reliability. G-Memory achieves better mean quality on several tasks, including study-adverse and study-outcome. Over reuse cutoffs 2--3 alone, \system\ meets all six quality margins; the first-cutoff failures remain in Figure~\ref{fig:official-quality}. Each method has one dependent sequence, so these are descriptive comparisons.}
{rel-trial结果揭示仍存的失败：\system\ 在九次尝试中得到六个有效方案，ACE、AWM与G-Memory各为七个。study-adverse MAE比无记忆高\OfAdverseRegression\%，未达到门槛；其首截点没有有效方案。site-success均值虽有改善，但三次尝试中仅一次有效，仅看均值会夸大可靠性。G-Memory在包括study-adverse与study-outcome的多项任务上均值更好。仅看复用截点2--3，\system\ 满足全部六项质量门槛；首截点失败仍保留在图~\ref{fig:official-quality}中。每种方法只有一条相互依赖的序列，因此这些是描述性比较。}
